\section*{Capítulo \ref{ch:b2:measurement-statistics} --- Estadística de la medida química}

\begin{solution}{exo:b2:measurement-statistics:1}
$\bar V = 62.41/5 = \qty{12.482}{mL}$; $\sum(V_i - \bar V)^2 = 0.00328$, $s = \sqrt{0.00328/4} \approx
\qty{0.029}{mL}$; $s/\sqrt5 \approx \qty{0.013}{mL}$.
\end{solution}

\begin{solution}{exo:b2:measurement-statistics:2}
$t = 2.776$ para 4 grados de libertad: $12.482 \pm 2.776 \times 0.0128$, es decir, $(12.48 \pm 0.04)$~mL.
\end{solution}

\begin{solution}{exo:b2:measurement-statistics:3}
Las incertidumbres relativas de la masa y del volumen son del 0.04~\% y del 0.06~\%; combinadas
cuadráticamente,
\[\frac{u(c)}{c} = \sqrt{\Bigl(\frac{0.0002}{0.5000}\Bigr)^2 + \Bigl(\frac{0.15}{250.0}\Bigr)^2}
\approx 0.07~\%.\]
\end{solution}

\begin{solution}{exo:b2:measurement-statistics:4}
$\bar x = 1.5$, $\bar y = 0.31$, $S_{xx} = 5$, $S_{xy} = 0.99$: $b = 0.198$, $a = 0.31 - 0.198 \times 1.5
= 0.013$.
\end{solution}

\begin{solution}{exo:b2:measurement-statistics:5}
$|10.12 - 10.00|/(0.08/\sqrt5) = 0.12/0.036 \approx 3.4 > 2.776$: la diferencia es significativa, y el
método está sesgado (en unos \qty{+0.12}{mg/L}).
\end{solution}

\begin{solution}{exo:b2:measurement-statistics:6}
$\mathrm{pH} = -\log_{10}c$, $\mathrm d\,\mathrm{pH}/\mathrm dc = -1/(c\ln10)$:
$u(\mathrm{pH}) = (u(c)/c)/\ln10 = 0.02/2.303 \approx 0.009$.
\end{solution}

\begin{solution}{exo:b2:measurement-statistics:7}
$x_{\mathrm{LOD}} = 3 \times 0.0012/0.0098 \approx \qty{0.37}{\micro g/L}$; $x_{\mathrm{LOQ}} = 10
\times 0.0012/0.0098 \approx \qty{1.2}{\micro g/L}$.
\end{solution}

\begin{solution}{exo:b2:measurement-statistics:8}
La recta da \qty{3.11}{mg/L} en las disoluciones medidas; la muestra se había diluido cinco veces:
$3.11 \times 50.0/10.0 \approx \qty{15.6}{mg/L}$.
\end{solution}

\begin{solution}{exo:b2:measurement-statistics:9}
\qty{0.05}{mg/L} es la \omterm{def:b2:measurement-statistics:validation}{repetibilidad}, y \qty{0.15}{mg/L}, la \omterm{def:b2:measurement-statistics:validation}{reproducibilidad}. Cada laboratorio tiene su
propio pequeño \omterm{def:b2:measurement-statistics:validation}{sesgo} (calibrado, patrones, instrumento); estos \omterm{def:b2:measurement-statistics:validation}{sesgos} difieren de un laboratorio a otro y se suman
a la dispersión, de modo que la \omterm{def:b2:measurement-statistics:validation}{reproducibilidad} es siempre la mayor.
\end{solution}

\begin{solution}{exo:b2:measurement-statistics:10}
$\partial S/\partial a = 0$ se escribe $\sum(y_i - a - bx_i) = 0$: los \omterm{def:b2:measurement-statistics:calibration}{residuos} suman cero. Dividiendo por
$n$: $\bar y = a + b\bar x$, de modo que $(\bar x, \bar y)$ está sobre la recta.
\end{solution}

\begin{solution}{exo:b2:measurement-statistics:11}
En una etapa:
\[\sqrt{(0.006/1.000)^2 + (0.08/100.0)^2} \approx 0.61~\%.\]
Una etapa de diez veces:
\[\sqrt{(0.02/10.00)^2 + (0.08/100.0)^2} \approx 0.22~\%,\]
y dos de ellas combinadas cuadráticamente dan $0.22\sqrt2 \approx 0.30~\%$. La vía en dos etapas es el doble de precisa: la pipeta pequeña domina la vía
en una etapa.
\end{solution}

\begin{solution}{exo:b2:measurement-statistics:12}
$z = 1.6/\sqrt{0.4^2 + 0.5^2} \approx 2.5 > 2$: no son compatibles; uno de los laboratorios (al menos) tiene un \omterm{def:b2:measurement-statistics:validation}{sesgo}.
Con incertidumbres expandidas ($k = 2$), el primero da $9.6 \pm 0.8$, que incluye 10; el segundo,
$11.2 \pm 1.0$, enteramente por encima de 10. Antes de cualquier veredicto sobre el agua hay que resolver la discrepancia,
por ejemplo con un material de referencia.
% ledger: who:lead
\end{solution}

\begin{solution}{pb:b2:measurement-statistics:1}
\textbf{1.} El ácido modifica la atomización y la respuesta; los patrones y las muestras deben tener la misma
matriz para que la pendiente sea aplicable.
\textbf{2.} $\bar x = 10$, $\bar y = 0.1004$, $S_{xx} = 250$, $S_{xy} = 2.445$: $b = 0.00978$ por
\unit{\micro g/L}, $a = 0.1004 - 0.0978 = 0.0026$.
\textbf{3.} $+0.0004$, $-0.0005$, $+0.0006$, $-0.0013$, $+0.0008$.
\textbf{4.} Los signos alternan, sin curvatura: la recta es adecuada en el intervalo 0--\qty{20}{\micro g/L}.
\textbf{5.} $s_{y/x} = \sqrt{3.1 \times 10^{-6}/3} \approx 0.00102$.
\textbf{6.} El blanco (ácido, agua, horno) da por sí mismo una pequeña absorbancia.
\textbf{7.} La sensibilidad: absorbancia por \unit{\micro g/L} de plomo.
\textbf{8.} $\bar y_0 = 0.1010$; $s = 0.0030$.
\textbf{9.} $x_0 = (0.1010 - 0.0026)/0.00978 \approx \qty{10.06}{\micro g/L}$.
\textbf{10.} $s_{x_0} = (0.00102/0.00978)\sqrt{1/3 + 1/5 + (0.1010 - 0.1004)^2/(0.00978^2 \times 250)}
\approx 0.104 \times 0.730 \approx \qty{0.076}{\micro g/L}$.
\textbf{11.} $n - 2 = 3$; $t = 3.182$.
\textbf{12.} $10.06 \pm 3.182 \times 0.076 = (10.06 \pm 0.24)$~\unit{\micro g/L}.
\textbf{13.} El término $1/m$ baja de 1 a $1/3$: $s_{x_0}$ disminuye en alrededor de un tercio.
\textbf{14.} $f = 50.50/50.0 = 1.010$; $10.06 \times 1.010 \approx \qty{10.16}{\micro g/L}$.
\textbf{15.} $f = 1 + V_2/V_1$: $u^2(f) = (u(V_2)/V_1)^2 + (V_2u(V_1)/V_1^2)^2 = (0.005/50.0)^2 +
(0.50 \times 0.03/2500)^2$, $u(f) \approx 1.0 \times 10^{-4}$, un 0.01~\% relativo.
\textbf{16.} No: 0.01~\% frente a $0.076/10.06 \approx 0.75~\%$.
\textbf{17.} $s_{\mathrm{blank}} \approx 0.00115$; $x_{\mathrm{LOD}} = 3 \times 0.00115/0.00978 \approx
\qty{0.35}{\micro g/L}$.
\textbf{18.} $10 \times 0.00115/0.00978 \approx \qty{1.2}{\micro g/L}$.
\textbf{19.} Sí, con mucho.
\textbf{20.} Sí: de 9.92 a \qty{10.40}{\micro g/L} en el agua del grifo.
\textbf{21.} $(0.104 - 0.0026)/0.00978 \times 1.010 \approx \qty{10.5}{\micro g/L}$ y, leída sola,
un ``por encima del valor guía'' categórico que los datos no respaldan.
\textbf{22.} La regla usada para construir el intervalo captura el valor verdadero en el 95~\% de las series a
las que se aplica; no es una probabilidad sobre este valor concreto.
\textbf{23.} Más lecturas replicadas y más patrones cerca de \qty{10}{\micro g/L} (los términos $1/m$ y
$1/n$), o un método más sensible; el intervalo solo se estrecha lentamente, como $1/\sqrt{\text{número}}$.
\textbf{24.} Analizando un agua de referencia certificada, o añadiendo a la muestra una cantidad conocida de
plomo y comprobando la recuperación.
\textbf{25.} Cerca de \qty{10}{\micro g/L} la incertidumbre de los métodos de rutina es de unos pocos por ciento, como aquí: un
valor guía más bajo no podría ser comprobado de forma fiable por los laboratorios ordinarios.
\textbf{26.} $\boldsymbol{(10.16 \pm 0.24)}$~\unit{\micro g/L} al 95~\%: el intervalo contiene
\qty{10}{\micro g/L}, de modo que la medida no demuestra ni que el agua supere el valor guía ni que lo cumpla
con margen.
\end{solution}
